% 核心观点：表示与反馈提供条件性的改进路径，实际任务收益须由匹配的闭环实验确认。
% 叙述逻辑：互补贡献与资源成本 -> 观测、覆盖和学习局限 -> 有边界的结论。
\section{Discussion and Conclusion}
\subsection{Discussion}
本文讨论的核心是有限机器人示范下的任务完成表现，
语义引导关系表示与历史条件反馈分别针对表示不足和执行信息不足。
两条路径应依据交叉对照独立评价；完整方法的提升不能替代对各环节的验证，
单一数据规模的结果也不能替代样本效率曲线。
离线语义资源与额外残差训练成本属于方法成本，须与性能一并报告。
当前三组front关系是给定操作任务的实例；新任务需要重新确定对象、关系和有效性约定，
并报告相应标注与训练成本。跨硬件结果不等同于跨任务通用性，无多任务证据时结论限定于已验证任务。
语义和几何监督提供任务相关归纳偏置，但不能补全未覆盖的接触动力学，
也不能恢复观测中无法识别的任意隐藏状态。
质量分数可能无法发现稳定的系统误分，历史模型也可能依赖错误相关性。
当前查询监督不等同于显式完成判定，监督残差亦不具备强化学习的策略改进保证。
分析中的覆盖、关系对任务后果的近似充分性和后果映射正则性是适用条件，
不是网络结构自动保证的性质。
当前方法没有学习任务后果函数；有限动作监督能否实现理论上的信息价值，
仍取决于实际决策学习与闭环验证。

式~\eqref{eq:successbound}以总体模仿风险$R_\nu(\pi)$为输入，
给出覆盖与敏感性假设下的任务失败概率上界，并非从有限示范上的训练损失直接得到的保证。
若要以经验误差替代总体风险，还需控制学得策略在相同动作目标、尺度和损失下的泛化差距，
并考虑轨迹内样本相关性以及利用数据进行模型选择的影响。
动作块训练损失与最终执行动作的平方误差也不一定相同，不能未经对应关系分析就直接代入该界。
本文不建立从有限样本训练误差到闭环任务成功率的端到端统计保证。

数据多样性、按完整episode划分数据以及保留不参与模型选择的测试集，
可用于降低采样局限与数据泄漏风险，并检验未见轨迹上的泛化表现；
但随机划分episode本身不保证统计独立，这些措施也不等同于形式化泛化界。
实验中应分别报告未见轨迹上的动作误差与闭环任务成功率，
区分同分布泛化和新位置、物体或环境带来的分布变化，
并说明分割器与策略各自的数据使用范围，避免将上游模型已使用的数据视为整个方法完全未见的测试数据。
\draftnote{待补：根据实际数据划分、模型选择协议及上游分割训练记录核实上述评价范围，再结合实验结果讨论泛化差距。}

\draftnote{待补讨论：依据最终失败案例说明方法在哪些条件下失效，跨平台和任务迁移的代价，以及自动标注的人工开销。}
\subsection{Conclusion}
本文围绕两个互补缺口构建研究框架：有限动作监督下任务关系可能未被充分保留，
而生成动作块时的信息可能不足以支持执行中的动态修正。
语义与几何监督用于构建任务关系表示，计划条件历史残差用于利用新增执行信息。
条件性分析区分任务相关信息损失、历史的潜在决策价值、实际决策误差与目标槽位时效；
这些依据解释可能的改善路径，不构成完整系统的无条件性能保证。
\draftnote{待补正文：基于最终验证结果，简洁总结问题、方法、获得支持的结论及剩余局限；暂不声明统计显著提升或通用收敛。}
